\chapter*{Preface}
Tome~5 ended with the prime number theorem, its error term, and the first sieves. It left the primes looking regular from far away and mysterious from close up. How long can a gap between consecutive primes be? Are there infinitely many pairs of primes at bounded distance? Is every large odd number a sum of three
primes? How evenly are the primes spread among the residue classes of a modulus that grows with $x$? Does the M\"obius function cancel on average for reasons that have nothing to do with the zeros of $\zeta$? This book proves the answers that are known, most of them found in the last hundred years and several in the last fifteen.

The first part is about zeros. Few zeros of $\zeta$ lie far from the critical line, and this is enough to find primes in every interval $\left]x,x+x^{0.76}\right]$. Siegel's theorem shows that a Dirichlet $L$-function cannot vanish too close to $1$, at the price of a constant that nobody can compute, and it gives the growth of class numbers promised in Tome~12.

The second part is about averages. Sieves meet the parity problem, which explains why they cannot find primes alone. The large sieve proves the Bombieri--Vinogradov theorem: on average over moduli up to $\sqrt x$, primes in progressions behave as the generalised Riemann hypothesis predicts. Vinogradov's method for exponential sums over primes and the circle method
of Hardy and Littlewood count partitions, sums of powers, and sums of three primes.

The third part reaches the frontier: the work of Goldston, Pintz and Y\i ld\i r\i m and of Maynard on small gaps between primes, Hal\'asz's theorem on multiplicative functions, the sum-product phenomenon, the cap set problem, Gowers norms, and Montgomery's pair correlation of the zeros of $\zeta$, where number theory meets random matrices.

The aim throughout is to understand why the theorems are true. Every result is proved, except a short list of deep theorems that are quoted, each marked as such where it appears, and never used silently. Pause and think boxes ask you to guess before reading on. Each chapter ends with exercises, graded from warm-ups ($\star$) to small
gems ($\star\star\star$), each with a hint and a complete solution that first explains how to find the idea and then writes it up. A companion Python notebook for each chapter checks the numbers in the text and opens a window on computer science: generating primes for cryptography, primality tests, sieving algorithms, pseudorandomness, fast
Fourier transforms, property testing and the fast evaluation of $\zeta$.

The book assumes Tome~5, and uses Tomes~4, 7, 10 and 12 where they help. Its beauty is in the way the primes, which no formula describes, obey laws that can be proved: deterministic, yet statistically as regular as a random sequence, and the chapters that follow try to show how much of that regularity can be reached.
